% Fichier généré automatiquement par tools/generate_corrections.py.
% Source : TEC/TEC-04-Meq-Jeq/25_Cheville/25_Cheville.tex
% Statut : partiel (3/5 question(s) renseignée(s), 0 reprise(s)).
\begin{corrigebox}[Corrigé partiel extrait de la banque]
\CorrigeQuestion{1}
D'après le diagramme de définition des blocs et le diagramme des exigences, les rapports de transmission doivent être : 
\begin{itemize}
\item pour l'axe de tangage : $\dfrac{N_{\text{moteur}}}{N_{\text{Tangage}}}=138,33$ au minimum; 
\item pour l'axe de roulis :  $\dfrac{N_{\text{moteur}}}{N_{\text{Roulis}}}= 197,61$ au minimum.
\end{itemize}
\CorrigeQuestion{2}
\CorrigeACompleter{Aucun élément de réponse n'est présent dans la source d'origine pour cette question.}
\CorrigeQuestion{3}
\CorrigeACompleter{Aucun élément de réponse n'est présent dans la source d'origine pour cette question.}
\CorrigeQuestion{4}
$$
R_T = (-1)^n \dfrac{80\cdot 47 \cdot 58 \cdot 36}{20\cdot 25\cdot 12 \cdot 10 } = 130,85
$$

Ceci est inférieur à ce qui est préconisé par le cahier des charges. 

Pour respecter le cahier des charges, on peut :
\begin{itemize}
\item choisir un autre moteur;
\item changer le nombre de dents d'une des roues. Il suffirait pour cela que,  par exemple, la roue de sortie comporte 39 dents. 
\end{itemize}
\CorrigeQuestion{5}
Le rapport de transmission du second train est de 201,3 ce qui est compatible avec le cahier des charges.
\end{corrigebox}
